\documentclass[11pt]{article}
\usepackage[margin=2.4cm]{geometry}
\usepackage{fontspec}
\setmainfont{DejaVu Sans}
\usepackage{juliamono}
\usepackage{fvextra}
\usepackage{xcolor}
\usepackage{xurl}
\usepackage{hyperref}
\usepackage{amssymb}
\newcommand{\checkbox}{$\square$}
\newcommand{\answerchoices}{\checkbox\ YES \quad \checkbox\ NO \quad \checkbox\ CANNOT DETERMINE}
\newcommand{\decisionexplanation}{%
  \par\smallskip\textbf{If NO, give the reason for the mismatch (required). If
  CANNOT DETERMINE, state what context or evidence is missing:}\par
  \rule{\linewidth}{0.4pt}\par\vspace{7pt}\rule{\linewidth}{0.4pt}\par}
\setlength{\parindent}{0pt}\setlength{\parskip}{5pt}
\setlength{\emergencystretch}{2em}
\DefineVerbatimEnvironment{Vsrc}{Verbatim}{fontsize=\small,frame=single,
  breaklines=true,breakanywhere=true}
\DefineVerbatimEnvironment{Vlean}{Verbatim}{fontsize=\scriptsize,frame=single,
  numbers=left,numbersep=5pt,breaklines=true,breakanywhere=true}
\hypersetup{colorlinks=true,urlcolor=blue}

\begin{document}
\begin{center}{\LARGE Proof Verification Document}\\[2pt]{\large mse\_2605396\_\_ans2605536}\\[2pt]{\small math.StackExchange mined pool (confirmed)}\end{center}
\begin{center}\fcolorbox{blue!60!black}{blue!4}{\parbox{0.92\linewidth}{\textbf{Make three independent decisions.} Follow the evidence chain: original source $\longrightarrow$ challenged claim $\longrightarrow$ proposed counterexample $\longrightarrow$ Lean theorem. A record is valid only when all three links are correct. Complete the human sign-off before reading the automated review appendix; those earlier results are background, not an answer key.}}\end{center}
\textbf{Use only this document, except for the public source link.} Everything needed for Questions 2 and 3 is included here. For Question 1, open the public source link when one is shown. If there is no public link, the complete stored problem and solution are included here. If the information is still not enough, choose CANNOT DETERMINE.
For each question, choose \textbf{YES}, \textbf{NO}, or \textbf{CANNOT DETERMINE}. Choose CANNOT DETERMINE if the document does not provide enough information for you to decide; do not mark YES if you are uncertain. If you find a specific mismatch, choose NO and explain it in the space below that question. Do not rewrite the claim, add assumptions, or substitute a different counterexample.
\textbf{How copied text is shown.} The problem, quoted passage, and proposed counterexample come from the stored record. The Lean theorem statement and Lean files are copied exactly from the certificate. Some long lines continue on the next printed line; this is only page formatting and does not change the copied text. Line numbers are added only to help you find passages. If the stored record contains a character that cannot be printed, the document shows \texttt{[U+XXXX CONTROL]} in its place.
\begin{description}
\item[Challenged claim.] The exact statement from the source that this record says is false.
\item[Proposed counterexample.] A concrete case that should satisfy the claim's assumptions but make its conclusion false.
\item[Lean certificate.] A Lean file proving a formal theorem. Successful compilation proves the theorem as written, not that it is the right translation of the source.
\end{description}
\section*{Question 1 --- Does this proof verification document accurately represent the source?}
\textbf{What to decide.} Compare the challenged passage below with the original source---either the public link or the complete source text in this document. Answer YES only if they make the same claim and include the same assumptions or conditions. If this document changes the meaning, leaves out a condition, or points to a different passage, answer NO.
\subsection*{Original source}
\textbf{Source record:}
\begin{Vsrc}
qedbench_certificate_availability_triage/data/records_mined_pilot.jsonl
\end{Vsrc}
\textbf{Public source:}\par
\url{https://math.stackexchange.com/a/2605536}\par
\subsection*{Problem or question from the source (verbatim)}
\begin{Vsrc}
Title: Proving that no two teams in a tournament win same number of games

Assume we've a tournament with $93$ teams. And between every $19$ teams, there is one team that has lost all of its $18$ games to other teams, and there is one team that has won all of its $18$ games against the other teams.
Between all teams, there is one match and only possible results for a match is win and lose.

My own idea is to try to give a proof by contradiction, by assuming two teams share same number of games won.
\end{Vsrc}
\subsection*{Challenged passage recorded from the source (stored text)}
\begin{Vsrc}
This gives you that there are teams with $k+1, k+2, \dots, n$ wins.
\end{Vsrc}
\subsection*{Surrounding source lines (stored text with added line labels)}
The labels refer to lines in the solution text stored with the source record; the source-line content is not transliterated.
\begin{Vsrc}
0001 | Definition top/bottom element: In a set $S$ a team is the top element if it has won against everyone else. It is the bottom if it has lost against everyone else.
0002 | 
0003 | Definition: Let's say that a set of teams is $k$ absolute if every subset of size $k$ of these teams has both a top and a bottom element.
0004 | 
0005 | To prove: You want to show your teams form a total order in terms of the win/lose relation. Here is a general proof that a total order exists for any $n$ teams which are $k$ absolute whenever $n \geq 2k > 4$. In your case, $n, k$ are $93, 19$ respectively, fitting the above inequality.
0006 | 
0007 | Proof: Suppose you have $n \geq 2k$ nodes. We can prove they form a total order as follows. Pull off the top node to get a $k$ absolute set of $n-1$ elements. Then pull off the top node again to get a $k$ absolute set of $n-2$ elements. Repeat until you're left with $k$ elements. You are allowed to do this until you get to $k$ by the "top" part of the top-bottom lemma. This gives you that there are teams with $k+1, k+2, \dots, n$ wins. Now, do the same for the least elements, again using the top-bottom lemma- this time the "bottom" part. You end up showing there are teams with $1, 2, \dots, k$ wins; you can stop at $k$ even if $n > 2k$ becausee you've already shown the wins from $k+1$ onwards. Together, then these results show that we have some team with $1$ win, another with $2$ etc. all the way up to $n$. This implies uniqueness by the pigeonhole principle because we have $n$ pigeons (teams) and $n$ holes (number of wins) and there exists a team for each number of wins.
0008 | 
0009 | Lemma (Top/Bottom): We show that every $k$ absolute set $U$ of $n$ teams with $n \geq k > 2$ elements has a top and a bottom element. We prove this as follows. Let $S$ be the largest subset of $U$ that has a top element. We know by $k$ absoluteness of $U$ that $|S| \geq k$. Now, suppose $S \neq U$, towards a contradiction. Let $t$ be the top element of $S$. Then we can choose some $u \in U$ s.t. $u > t$. Now, we can show that $u$ is the top element in $S \cup \{u\}$. To show this, consider any subset of $S$ of size $|S| - 1$ that also contains $t$. Call it $S'$. Because $k > 2$, every element in $S$ belongs to some such $S'$; a crucial point made by commenter Steve Kass is that this is not the case for $k = 2$. Now, by $k$ absoluteness, there must be a top element of $S' \cup \{u\}$. Well, because $t$ is greater than all elements existing other than $u$, none of them can be the top. Nor can $t$ be the top since $u>t$. So $u$ is the top element. So, $u$ is greater than all elements of $S'$. Since we've already mentioned that every element of $S$ is in some $S'$, because $k>2$, we know then that $u > s$ for every $s \in S$. Thus the set $S \cup \{u\}$ of size strictly greater than $S$, has a top element $u$, contradicting our assumption that $S$ was the top. Thus, $U$ must have a top element. The exact same argument can be applied for a bottom element.
\end{Vsrc}
\textbf{A YES means:} the proof verification document faithfully represents both the source context and the challenged claim, without dropping a condition or changing its meaning.
\section*{Question 2 --- Does the proposed counterexample disprove that claim?}
\textbf{Claim to test (repeated for convenience).}
\begin{Vsrc}
This gives you that there are teams with $k+1, k+2, \dots, n$ wins.
\end{Vsrc}
\subsection*{How the record says the argument uses this claim}
This is the record's description, not an established fact; check it against the source.
\begin{Vsrc}
The proof's enumeration step asserts that peeling off top elements yields teams with k+1, k+2, ..., n wins (and, with the bottom peel, teams with every win-count 1, 2, ..., n), so that the pigeonhole is applied to n teams and the n "holes" {1,...,n}. This win-count set is infeasible in an n-team round-robin tournament: each team plays only n-1 games, so no team can have n wins, and the multiset {1,...,n} has total weight n(n+1)/2, not the C(n,2)=n(n-1)/2 games actually played. The correct sets produced by the argument are {k,...,n-1} (top peel) and {0,...,n-k-1} (bottom peel); every index in the proof is shifted by +1.
\end{Vsrc}
\subsection*{Proposed counterexample (exact values or construction)}
\begin{Vsrc}
Witness: the 5-team transitive round-robin tournament (n = 5). All C(5,2) = 10 games: 0 beats 1,2,3,4; 1 beats 2,3,4; 2 beats 3,4; 3 beats 4. Win-vector (team 0..4) = (4,3,2,1,0), i.e. the win-counts are exactly the 5 distinct values {0,1,2,3,4}. This is precisely the "peeling" configuration the proof describes (every team has a different number of wins), so it cannot be dismissed as degenerate.

Take k = 2. The proof's sentence asserts "there are teams with k+1, k+2, ..., n wins", i.e. teams with 3, 4 and 5 wins. Recomputed on the witness: the win-counts actually present are {0,1,2,3,4}; the maximum win-count over all 5 teams is 4. Claim says some team has 5 wins; here max wins = 4 < 5, so no team has 5 wins. Exact reason: each team plays exactly n-1 = 4 games, so every win-count lies in {0,...,4} and the value n = 5 is out of range.

Counting check on the pigeonhole holes {1,...,n}: the claimed hole-set has total weight 1+2+3+4+5 = 15, while the tournament distributes exactly C(5,2) = 10 wins. Claim says total = 15, actual total = 10 (15 != 10).

Correct sets on the witness: top peel = {k,...,n-1} = {2,3,4}, bottom peel = {0,...,n-k-1} = {0,1,2} — every index in the proof's enumeration is shifted by +1, exactly as the target claim states.

Beyond this single witness, the script exhaustively enumerated every orientation of K_n for n = 2..7 (2 + 8 + 64 + 1024 + 32768 + 2097152 = 2,130,018 tournaments) and confirmed, with exact integer arithmetic: max win-count over ALL tournaments is always exactly n-1; no tournament ever has a team with n wins; no tournament ever has win-count multiset {1,...,n}; total wins is always C(n,2); and for no n and no k in 0..n-1 does any tournament realize the asserted top-peel set {k+1,...,n}. So every round-robin tournament is a counterexample to the enumeration step. Script runtime ~24 s.
\end{Vsrc}
\subsection*{Why the record says the counterexample works}
Check this explanation independently; it is evidence to assess, not an answer key.
\begin{Vsrc}
The asserted statement is finite and combinatorial: in a round-robin tournament on n teams, the peeling step produces teams whose win-counts are k+1, k+2, ..., n (and, with the bottom peel, the n teams occupy the n holes {1,...,n}). Its refutable content is (A) some team has n wins, (B) the win-count multiset can be {1,...,n}, and (C) the implied total of wins is n(n+1)/2.

All three fail on every tournament. In a round-robin each team plays exactly n-1 games, so its win-count is an integer in {0,...,n-1}; the value n is unattainable, hence the top-peel set {k+1,...,n} is never realized for any k, and the multiset {1,...,n} is never realized. The counting check is independent confirmation: the tournament plays exactly C(n,2) = n(n-1)/2 games and therefore distributes exactly n(n-1)/2 wins, whereas the claimed hole-set {1,...,n} carries weight n(n+1)/2, an excess of exactly n.

The script verifies this not by appeal to the argument but by brute force: it enumerates every one of the 2^C(n,2) orientations of K_n for n = 2 through 7 (2,130,018 tournaments in total), computes each exact integer win-vector, and records the maximum win-count, whether any team ever reaches n wins, whether the multiset {1,...,n} ever occurs, whether the total ever differs from C(n,2), and whether the asserted set {k+1,...,n} is ever realized for any k. Results: max win-count is exactly n-1 for every n; zero tournaments with a team having n wins; zero tournaments with multiset {1,...,n}; total wins always C(n,2); no k admissible for any n. The one shortcut used (skipping the per-k subset test when max win-count < n) is exact, not heuristic, since {k+1,...,n} always contains n.

All arithmetic is Python int; itertools.product/combinations drive a deterministic enumeration with no randomness and no floating point in any decision.

Honest caveat: I was given only the claim and the verbatim sentence, not the source problem, so I fixed n = number of teams as the target claim specifies. If the underlying source were phrased with n+1 teams (or with "wins" meaning something other than games won), the indexing critique would need to be re-checked against that convention; under the stated convention the refutation is airtight and in fact universal.
\end{Vsrc}
\subsection*{How to decide}
\begin{enumerate}
\item List every assumption that applies to the challenged claim.
\item Check that the proposed example satisfies every one of those assumptions.
\item Recompute the claimed outcome and the actual outcome using exact reasoning.
\item Confirm that the conflict is with the challenged claim itself, not an unrelated sentence.
\end{enumerate}
\textbf{A YES means:} the example is fully specified, satisfies all relevant assumptions, and makes the challenged claim false. If a necessary value, condition, or calculation is missing, choose CANNOT DETERMINE.
\section*{Question 3 --- Does the final Lean theorem statement verify this same counterexample?}
\textbf{What to decide.} A Lean file can compile while proving the wrong theorem. Compare these three links independently:
\begin{enumerate}
\item The Lean definitions represent the same mathematical objects, assumptions, and operations as the challenged claim.
\item The Lean construction uses the same concrete values or example shown in Question 2.
\item The final theorem statement asserts that the challenged claim fails on that example, not merely a weaker or nearby statement.
\end{enumerate}
\subsection*{Selected final Lean theorem statement (verbatim UTF-8 excerpt)}
This is an exact contiguous slice of the certificate file, from the theorem keyword through the original \texttt{:=} token where its proof begins. It is shown first so the formal target is not buried in the complete file.
\begin{Vlean}
theorem mse_2605396__ans2605536_witness5_refutes :
    -- (A) the asserted top peel for k = 2 is not realized; the shifted-down one is
    (natListEq (claimedTopPeel 2 5) [3, 4, 5] = true)
  ∧ (natListEq (correctTopPeel 2 5) [2, 3, 4] = true)
  ∧ (realizesSet (winVec 5 (transTour 5)) (claimedTopPeel 2 5) = false)
  ∧ (realizesSet (winVec 5 (transTour 5)) (correctTopPeel 2 5) = true)
  ∧ (hasWinCount (winVec 5 (transTour 5)) 5 = false)
  ∧ (hasWinCount (winVec 5 (transTour 5)) 4 = true)
  ∧ (natMax (winVec 5 (transTour 5)) = 4)
    -- (B) the holes are {0,…,4}, not {1,…,5}
  ∧ (natListEq (sortNat (winVec 5 (transTour 5))) (claimedHoles 5) = false)
  ∧ (natListEq (sortNat (winVec 5 (transTour 5))) (actualHoles 5) = true)
  ∧ (realizesSet (winVec 5 (transTour 5)) (claimedHoles 5) = false)
  ∧ (realizesSet (winVec 5 (transTour 5)) (actualHoles 5) = true)
    -- (C) the tally: 15 claimed, 10 actual
  ∧ (claimedTotalWins 5 = 15)
  ∧ (actualTotalWins 5 = 10)
  ∧ (natSum (winVec 5 (transTour 5)) = 10) :=
\end{Vlean}
\subsection*{Project-local imported Lean files (complete, verbatim UTF-8)}
This section follows the certificate's import commands recursively and prints every repository-owned Lean source file on which the final theorem statement depends. Each file is copied in full, including its definitions and controls. External Mathlib, Std, and Lean Init source is not duplicated.
No project-local imported Lean files were found. External Mathlib, Std, and Lean Init source is not duplicated in this proof verification document.
\subsection*{Complete counterexample Lean file (verbatim UTF-8)}
The complete Lean file is included below so you can check every definition and supporting result used by the final theorem statement. Use it together with the project-local imported Lean files above.
\textbf{Lean certificate file:}
\begin{Vsrc}
partB_feasibility/mined_certificates/mse_2605396__ans2605536.lean
\end{Vsrc}
\begin{Vlean}
/-!
# `mse_2605396__ans2605536` — the peeling step asserts teams with `k+1, …, n` wins;
  no tournament on `n` teams ever has a team with `n` wins

* **Record id:** `mse_2605396__ans2605536`
* **Question:** <https://math.stackexchange.com/questions/2605396>
  ("Proving that no two teams in a tournament win same number of games")
* **Answer (the flagged solution):** <https://math.stackexchange.com/a/2605536> (accepted)
* **Route:** `SELF_CONTAINED_LEAN`.  Chosen because the shared claim DSL
  `qedbench_certificate_availability_triage/lean_certificates/QEDTemplateV2.lean` has **no**
  constructor over tournaments / orientations of `K_n` and no notion of a win-count vector:
  its witness types are `BankGraph` (undirected simple graph), `WDG` (weighted digraph used
  only for shortest-path / MST claims) and auxiliary `Nat` lists, and grepping the whole
  `ClaimV2` constructor list turns up nothing about tournaments, out-degrees, win counts,
  top/bottom elements or `k`-absoluteness (the nearest neighbours are `degSeqEq` and
  `stronglyRegular`, both undirected).  There is therefore nothing to instantiate; this file
  stands alone.  It has **no `import` line at all** (core Lean 4 `Init` only), imports nothing
  from the trusted template, and modifies nothing outside itself.
* **Text licence:** CC BY-SA (Stack Exchange); the sentences below are quoted for
  identification.

## Verbatim claim being refuted

Line `0007` of the line-numbered solution text (the "Proof" paragraph).  The refuted
assertions are the two enumeration sentences and the pigeonhole tally that rests on them:

> "Repeat until you're left with $k$ elements. You are allowed to do this until you get to $k$
> by the "top" part of the top-bottom lemma. **This gives you that there are teams with
> $k+1, k+2, \dots, n$ wins.** Now, do the same for the least elements … You end up showing
> there are teams with $1, 2, \dots, k$ wins … **Together, then these results show that we have
> some team with $1$ win, another with $2$ etc. all the way up to $n$.** This implies uniqueness
> by the pigeonhole principle because we have $n$ pigeons (teams) and $n$ holes (number of
> wins) and there exists a team for each number of wins."

Its refutable arithmetic content, exactly as isolated by the search script
`partB_feasibility/counterexample_search_mined/mse_2605396__ans2605536.py`:

* **(A)** some team has `n` wins (every asserted top-peel set `{k+1, …, n}` contains `n`);
* **(B)** the `n` teams occupy the `n` holes `{1, 2, …, n}`, i.e. the win-count multiset is
  `{1, …, n}`;
* **(C)** hence the total number of wins is `1 + 2 + ⋯ + n = n(n+1)/2`.

All three are false for **every** tournament on `n` teams: each team plays exactly `n − 1`
games, so every win-count lies in `{0, …, n−1}`, the value `n` is unattainable, and the
tournament distributes exactly `C(n,2) = n(n−1)/2` wins, short of `n(n+1)/2` by exactly `n`.
The correct enumeration is shifted down by one: the top peel yields `{k, …, n−1}` and the
holes are `{0, 1, …, n−1}`.

### SCOPE — read before citing this file

Refuted here is the **enumeration/indexing step** of line `0007` only.  The answer's headline
conclusion — that in such a tournament no two teams win the same number of games — is
**true, and is positively certified below** (Part 4, `…_conclusion_of_the_answer_is_true`):
on both witnesses the win counts are pairwise distinct.  The Top/Bottom lemma of line `0009`
is untouched.  What fails is the arithmetic the pigeonhole argument is run on: the argument as
written needs a team with `n` wins, which no tournament has, and counts `n(n+1)/2` wins where
only `n(n−1)/2` exist.  This is an internal off-by-one in an otherwise sound proof sketch, not
a refutation of the theorem.

## Witnesses

1. **The triage witness** (`witness_description`): the transitive round-robin tournament on
   `n = 5` teams — `0` beats `1,2,3,4`; `1` beats `2,3,4`; `2` beats `3,4`; `3` beats `4`.
   Win vector `(4,3,2,1,0)`: five *distinct* win counts, so it cannot be dismissed as a
   degenerate configuration — it is exactly the "peeling" shape the proof describes.  With
   `k = 2` the proof asserts teams with `3, 4, 5` wins; the win counts present are
   `{0,1,2,3,4}`, so no team has `5`.  Total wins `10 = C(5,2)`, against the claimed
   `1+2+3+4+5 = 15`.
2. **An in-regime witness.**  The proof's standing hypothesis is `n ≥ 2k > 4`, which
   `n = 5, k = 2` does **not** satisfy (`2k = 4` is not `> 4`).  So the file also certifies the
   transitive tournament on `n = 6` with `k = 3`, where `6 ≥ 6 > 4` holds *and* the tournament
   is verified `3`-absolute (every 3-subset has a top and a bottom element) — i.e. every
   hypothesis of line `0007` is met and its conclusion still fails: no team has `6` wins,
   the asserted top peel `{4,5,6}` is not realized, total wins `15 ≠ 21`.

## Modelling judgments (the kernel checks the arithmetic, not the faithfulness of these)

1. **Tournament = orientation of `K_n`.**  `pairsLt n` lists the pairs `(i,j)`, `i < j`, in a
   fixed order; a tournament is a `List Bool` of that same length, entry `true` meaning "`i`
   beat `j`", `false` meaning "`j` beat `i`".  Every such list is a legal round-robin
   tournament and every round-robin tournament is such a list, so the encoding is exactly the
   set of tournaments — no well-formedness side condition can silently exclude a
   counterexample.  Ties are impossible by construction, matching the question's "only
   possible results for a match is win and lose".
2. **Win count = number of pairs whose winner is that team** (`winCount`), i.e. out-degree.
   `winVec n o` is the vector of win counts of teams `0, …, n−1`, in that order.
3. **Teams are 0-indexed**; the *win counts* are plain naturals, so the comparison with the
   proof's `1, …, n` is index-convention-free: the assertion "a team with `n` wins" is about
   a number of games won, not about a team label.
4. **The bottom-peel sentence is NOT claimed to fail.**  `{1, …, k}` is realizable (both
   witnesses do have teams with `1` and `2` wins), and this file makes no claim against it.
   What fails is the *top* peel `{k+1, …, n}` — for every `k` and every tournament — and
   therefore the union `{1, …, n}` and the pigeonhole tally.  The certificate is deliberately
   narrow on this point.
5. **`k`-absoluteness is modelled literally** from line `0003`: `kAbsolute n o k` says every
   `k`-element subset of the teams has a top element (beats all others in the subset) and a
   bottom element (loses to all others in the subset).  It is *discriminating*, not vacuously
   true: the 3-cycle tournament on 3 teams is shown to be non-`3`-absolute (Part 5, NC3).
6. **Exhaustive sweep budget.**  The kernel sweep covers every orientation of `K_n` for
   `n = 2,3,4,5` (2 + 8 + 64 + 1024 = 1098 tournaments, fully enumerated, no sampling).  The
   record's Python script additionally covered `n = 6` and `n = 7` (2,130,018 tournaments in
   total); those two sizes are **not** re-derived here — kernel budget only, not because they
   behave differently (the `n = 6` transitive witness *is* checked individually).
7. **Set-membership reading of "there are teams with `v₁, …, v_r` wins"**: `realizesSet`
   requires, for each listed value, *some* team with exactly that many wins.  This is the
   weakest reasonable reading — it does not demand distinct teams or an exact multiset — so
   refuting it refutes every stronger reading too.
8. Everything is discharged by kernel `decide` on closed `Bool`-valued terms.  No
   `native_decide`, no `sorry`, no axioms, no imports.
-/

set_option maxRecDepth 100000
set_option maxHeartbeats 4000000

namespace mse_2605396__ans2605536

/-! ## Part 0. Encoding: tournaments as orientations of `K_n` -/

/-- All pairs `(i,j)` with `i < j < n`, in lexicographic order.  This is the game schedule of
a round-robin tournament on `n` teams; its length is `C(n,2)`. -/
def pairsLt (n : Nat) : List (Nat × Nat) :=
  (List.range n).foldr (fun i acc =>
    ((List.range n).filter (fun j => decide (i < j))).foldr
      (fun j acc2 => (i, j) :: acc2) acc) []

/-- The winner of each scheduled game: entry `true` at pair `(i,j)` means `i` beat `j`. -/
def winnersOf : List (Nat × Nat) → List Bool → List Nat
  | [], _ => []
  | (i, j) :: ps, bs =>
      match bs with
      | [] => []
      | b :: bs' => (if b then i else j) :: winnersOf ps bs'

/-- How many games team `t` won, given the winner list. -/
def winCount (ws : List Nat) (t : Nat) : Nat :=
  (ws.filter (fun w => w == t)).length

/-- The win vector of teams `0, …, n−1` from an already-computed winner list. -/
def winVecOf (n : Nat) (ws : List Nat) : List Nat :=
  (List.range n).map (fun t => winCount ws t)

/-- The win vector of the tournament `o` on `n` teams. -/
def winVec (n : Nat) (o : List Bool) : List Nat :=
  winVecOf n (winnersOf (pairsLt n) o)

/-- The transitive tournament on `n` teams: `i` beats `j` whenever `i < j`. -/
def transTour (n : Nat) : List Bool := (pairsLt n).map (fun _ => true)

/-! ### Small list utilities (kept first-order so every statement is a closed `Bool` term) -/

def natSum : List Nat → Nat
  | [] => 0
  | a :: as => a + natSum as

def natMax : List Nat → Nat
  | [] => 0
  | a :: as => let m := natMax as; if a ≤ m then m else a

def insSorted (x : Nat) : List Nat → List Nat
  | [] => [x]
  | a :: as => if x ≤ a then x :: a :: as else a :: insSorted x as

def sortNat : List Nat → List Nat
  | [] => []
  | a :: as => insSorted a (sortNat as)

/-- Equality of `Nat` lists, as a `Bool` function (no `DecidableEq` instance trusted). -/
def natListEq : List Nat → List Nat → Bool
  | [], [] => true
  | a :: as, b :: bs => (a == b) && natListEq as bs
  | _, _ => false

/-! ## Part 1. The claim, transcribed

`hasWinCount wv v` — some team has exactly `v` wins.
`realizesSet wv vs` — for every value in `vs` there is a team with that many wins:
the literal content of "there are teams with `v₁, v₂, …` wins" (modelling judgment 7). -/

def hasWinCount (wv : List Nat) (v : Nat) : Bool := wv.any (fun w => w == v)

def realizesSet (wv : List Nat) (vs : List Nat) : Bool := vs.all (fun v => hasWinCount wv v)

/-- The set the proof asserts the top peel produces: `{k+1, k+2, …, n}`. -/
def claimedTopPeel (k n : Nat) : List Nat := (List.range (n - k)).map (fun t => k + 1 + t)

/-- What the top peel actually produces: `{k, k+1, …, n−1}` (everything shifted down by one). -/
def correctTopPeel (k n : Nat) : List Nat := (List.range (n - k)).map (fun t => k + t)

/-- The proof's pigeonhole holes: `{1, 2, …, n}`. -/
def claimedHoles (n : Nat) : List Nat := (List.range n).map (fun t => t + 1)

/-- The holes that actually exist: `{0, 1, …, n−1}`. -/
def actualHoles (n : Nat) : List Nat := List.range n

/-- The total the proof's tally implies: `1 + 2 + ⋯ + n`. -/
def claimedTotalWins (n : Nat) : Nat := natSum (claimedHoles n)

/-- The number of wins actually distributed: one per game, i.e. `C(n,2)`. -/
def actualTotalWins (n : Nat) : Nat := (pairsLt n).length

/-! ## Part 2. Witness 1 — the transitive tournament on 5 teams (the triage witness) -/

/-- The witness's schedule and win vector are what the header says: 10 games, win vector
`(4,3,2,1,0)` — five distinct win counts, the "peeling" configuration itself. -/
theorem mse_2605396__ans2605536_witness5_shape :
    ((pairsLt 5).length = 10)
  ∧ ((transTour 5).length = 10)
  ∧ (natListEq (winVec 5 (transTour 5)) [4, 3, 2, 1, 0] = true)
  ∧ (natSum (winVec 5 (transTour 5)) = 10)
  ∧ (actualTotalWins 5 = 10) := by decide

/-- **Main certificate (witness 1).**  On the 5-team transitive tournament, with `k = 2`:

* (A) no team has `n = 5` wins — the maximum win count is `4`, so the asserted top peel
  `{3,4,5}` is *not* realized, while the true top peel `{2,3,4}` is;
* (B) the win-count multiset is `{0,1,2,3,4}`, not the asserted `{1,2,3,4,5}`;
* (C) the tournament distributes `10` wins, not the `1+2+3+4+5 = 15` the pigeonhole tally
  assumes — a shortfall of exactly `n = 5`. -/
theorem mse_2605396__ans2605536_witness5_refutes :
    -- (A) the asserted top peel for k = 2 is not realized; the shifted-down one is
    (natListEq (claimedTopPeel 2 5) [3, 4, 5] = true)
  ∧ (natListEq (correctTopPeel 2 5) [2, 3, 4] = true)
  ∧ (realizesSet (winVec 5 (transTour 5)) (claimedTopPeel 2 5) = false)
  ∧ (realizesSet (winVec 5 (transTour 5)) (correctTopPeel 2 5) = true)
  ∧ (hasWinCount (winVec 5 (transTour 5)) 5 = false)
  ∧ (hasWinCount (winVec 5 (transTour 5)) 4 = true)
  ∧ (natMax (winVec 5 (transTour 5)) = 4)
    -- (B) the holes are {0,…,4}, not {1,…,5}
  ∧ (natListEq (sortNat (winVec 5 (transTour 5))) (claimedHoles 5) = false)
  ∧ (natListEq (sortNat (winVec 5 (transTour 5))) (actualHoles 5) = true)
  ∧ (realizesSet (winVec 5 (transTour 5)) (claimedHoles 5) = false)
  ∧ (realizesSet (winVec 5 (transTour 5)) (actualHoles 5) = true)
    -- (C) the tally: 15 claimed, 10 actual
  ∧ (claimedTotalWins 5 = 15)
  ∧ (actualTotalWins 5 = 10)
  ∧ (natSum (winVec 5 (transTour 5)) = 10) := by decide

/-- The failure is not special to `k = 2`: for **every** `k ∈ {0,…,4}` the asserted top peel
`{k+1,…,5}` fails on this tournament, and for every such `k` the shifted-down set `{k,…,4}`
succeeds.  So no choice of `k` rescues the sentence. -/
theorem mse_2605396__ans2605536_witness5_every_k :
    ((List.range 5).all (fun k =>
        !realizesSet (winVec 5 (transTour 5)) (claimedTopPeel k 5)) = true)
  ∧ ((List.range 5).all (fun k =>
        realizesSet (winVec 5 (transTour 5)) (correctTopPeel k 5)) = true) := by decide

/-! ## Part 3. Witness 2 — inside the proof's own regime `n ≥ 2k > 4` (`n = 6`, `k = 3`)

Witness 1 sits at `n = 5, k = 2`, where `2k = 4` is not `> 4`; a reader could object that the
proof never claimed anything there.  This part removes the objection: the transitive
tournament on 6 teams is verified `3`-absolute, `6 ≥ 2·3 > 4` holds, and the conclusion of
line `0007` still fails. -/

/-- Did `i` beat `j` in tournament `o`?  Scans the schedule; `false` if `i`, `j` never met
(only when `i = j` or an index is out of range). -/
def arcLookup : List (Nat × Nat) → List Bool → Nat → Nat → Bool
  | [], _, _, _ => false
  | (a, b) :: ps, bs, i, j =>
      match bs with
      | [] => false
      | w :: bs' =>
          if (a == i && b == j) then w
          else if (a == j && b == i) then !w
          else arcLookup ps bs' i j

def beatsB (n : Nat) (o : List Bool) (i j : Nat) : Bool := arcLookup (pairsLt n) o i j

/-- All `k`-element sublists of `l` (order-preserving; structural in `l`). -/
def subsetsOfSize : Nat → List Nat → List (List Nat)
  | 0, _ => [[]]
  | _ + 1, [] => []
  | k + 1, x :: xs => ((subsetsOfSize k xs).map (fun s => x :: s)) ++ subsetsOfSize (k + 1) xs

/-- Line `0001`: `S` has a *top* element — a member beating every other member. -/
def hasTop (n : Nat) (o : List Bool) (S : List Nat) : Bool :=
  S.any (fun t => S.all (fun u => (u == t) || beatsB n o t u))

/-- Line `0001`: `S` has a *bottom* element — a member losing to every other member. -/
def hasBot (n : Nat) (o : List Bool) (S : List Nat) : Bool :=
  S.any (fun t => S.all (fun u => (u == t) || beatsB n o u t))

/-- Line `0003`: the team set is `k`-absolute — every `k`-subset has a top and a bottom. -/
def kAbsolute (n : Nat) (o : List Bool) (k : Nat) : Bool :=
  (subsetsOfSize k (List.range n)).all (fun S => hasTop n o S && hasBot n o S)

/-- The 6-team transitive tournament satisfies every hypothesis of line `0007` for `k = 3`:
it is `3`-absolute (and `4`- and `5`-absolute), and `n ≥ 2k > 4` holds numerically. -/
theorem mse_2605396__ans2605536_witness6_hypotheses_hold :
    (kAbsolute 6 (transTour 6) 3 = true)
  ∧ (kAbsolute 6 (transTour 6) 4 = true)
  ∧ (kAbsolute 6 (transTour 6) 5 = true)
  ∧ (6 ≥ 2 * 3) ∧ (2 * 3 > 4)
  ∧ ((pairsLt 6).length = 15)
  ∧ (natListEq (winVec 6 (transTour 6)) [5, 4, 3, 2, 1, 0] = true) := by decide

/-- **Main certificate (witness 2).**  With every hypothesis of line `0007` satisfied
(`n = 6 ≥ 2k = 6 > 4`, the set is `3`-absolute), its conclusion is false: no team has `6`
wins, the asserted top peel `{4,5,6}` is not realized, the holes are `{0,…,5}` not `{1,…,6}`,
and the tournament distributes `15` wins against the claimed `21`. -/
theorem mse_2605396__ans2605536_witness6_refutes :
    (natListEq (claimedTopPeel 3 6) [4, 5, 6] = true)
  ∧ (natListEq (correctTopPeel 3 6) [3, 4, 5] = true)
  ∧ (realizesSet (winVec 6 (transTour 6)) (claimedTopPeel 3 6) = false)
  ∧ (realizesSet (winVec 6 (transTour 6)) (correctTopPeel 3 6) = true)
  ∧ (hasWinCount (winVec 6 (transTour 6)) 6 = false)
  ∧ (natMax (winVec 6 (transTour 6)) = 5)
  ∧ (natListEq (sortNat (winVec 6 (transTour 6))) (claimedHoles 6) = false)
  ∧ (natListEq (sortNat (winVec 6 (transTour 6))) (actualHoles 6) = true)
  ∧ (realizesSet (winVec 6 (transTour 6)) (claimedHoles 6) = false)
  ∧ (claimedTotalWins 6 = 21)
  ∧ (actualTotalWins 6 = 15)
  ∧ (natSum (winVec 6 (transTour 6)) = 15)
  ∧ ((List.range 6).all (fun k =>
        !realizesSet (winVec 6 (transTour 6)) (claimedTopPeel k 6)) = true) := by decide

/-! ## Part 4. Exhaustive sweep: EVERY tournament refutes the enumeration step -/

/-- All `2^m` `Bool` lists of length `m` — i.e. all orientations of a schedule of `m` games. -/
def allBoolLists : Nat → List (List Bool)
  | 0 => [[]]
  | m + 1 => (allBoolLists m).foldr (fun l acc => (true :: l) :: (false :: l) :: acc) []

/-- The five checks, on a win vector: no team with `n` wins; every win count `< n`; the
multiset is not `{1,…,n}`; the total is exactly `C(n,2)`; and no `k` realizes `{k+1,…,n}`. -/
def sweepOkWv (n : Nat) (wv : List Nat) : Bool :=
  !hasWinCount wv n
  && decide (natMax wv < n)
  && !natListEq (sortNat wv) (claimedHoles n)
  && (natSum wv == actualTotalWins n)
  && (List.range n).all (fun k => !realizesSet wv (claimedTopPeel k n))

def sweepOk (n : Nat) (o : List Bool) : Bool := sweepOkWv n (winVec n o)

/-- Every orientation of `K_n` passes all five checks. -/
def allSweepOk (n : Nat) : Bool :=
  (allBoolLists (pairsLt n).length).all (fun o => sweepOk n o)

/-- **Exhaustive certificate.**  Over *all* 2 + 8 + 64 + 1024 = 1098 tournaments on
`n = 2,3,4,5` teams, without exception: no team ever has `n` wins, no win-count multiset is
ever `{1,…,n}`, the total is always `C(n,2)`, and the asserted top-peel set `{k+1,…,n}` is
realized for no `k` whatsoever.  The refutation is therefore universal, not a lucky witness:
*every* round-robin tournament refutes the enumeration step. -/
theorem mse_2605396__ans2605536_exhaustive_sweep :
    (allSweepOk 2 = true)
  ∧ (allSweepOk 3 = true)
  ∧ (allSweepOk 4 = true)
  ∧ (allSweepOk 5 = true) := by decide

/-- Non-vacuity of the sweep: the enumerated families really have `2^C(n,2)` members, so the
theorem above is not true merely because the quantified list is empty or truncated. -/
theorem mse_2605396__ans2605536_sweep_nonvacuous :
    ((allBoolLists (pairsLt 2).length).length = 2)
  ∧ ((allBoolLists (pairsLt 3).length).length = 8)
  ∧ ((allBoolLists (pairsLt 4).length).length = 64)
  ∧ ((allBoolLists (pairsLt 5).length).length = 1024)
  ∧ ((allBoolLists 3).all (fun l => l.length == 3) = true) := by decide

/-- **The answer's actual conclusion is TRUE on both witnesses** — this file attacks the
enumeration step, not the theorem.  On each witness the win counts are pairwise distinct
(`sorted win vector = {0,…,n−1}`, which has no repeats), exactly as the question asks. -/
theorem mse_2605396__ans2605536_conclusion_of_the_answer_is_true :
    (natListEq (sortNat (winVec 5 (transTour 5))) (actualHoles 5) = true)
  ∧ (natListEq (sortNat (winVec 6 (transTour 6))) (actualHoles 6) = true)
  ∧ ((List.range 5).all (fun v => winCount (winVec 5 (transTour 5)) v == 1) = true)
  ∧ ((List.range 6).all (fun v => winCount (winVec 6 (transTour 6)) v == 1) = true) := by decide

/-! ## Part 5. NEGATIVE CONTROLS

Each conjunct below would fail to compile if the modelling were vacuous or rigged.  Together
they rule out the ways this certificate could be hollow: a `realizesSet` that never returns
`true`, a `winVec` that is constant or ignores the orientation, a `kAbsolute` that is
trivially true, and an encoding that cannot express a team with `n` wins at all. -/

/-- (NC1) `realizesSet` is **not** constantly false: on the very same witness, with the very
same evaluator, the *repaired* claim (every index shifted down by one) comes out `true` — for
the top peel at every `k`, for the full hole set, and for each individual value `0,…,4`.  The
refutation is therefore located precisely at the `+1`, not in a broken predicate. -/
theorem mse_2605396__ans2605536_nc1_repaired_claim_succeeds :
    (realizesSet (winVec 5 (transTour 5)) (actualHoles 5) = true)
  ∧ (realizesSet (winVec 5 (transTour 5)) [0, 1, 2, 3, 4] = true)
  ∧ ((List.range 5).all (fun v => hasWinCount (winVec 5 (transTour 5)) v) = true)
  ∧ (realizesSet (winVec 6 (transTour 6)) (actualHoles 6) = true)
  ∧ (realizesSet (winVec 5 (transTour 5)) [] = true) := by decide

/-- The reversed tournament: `j` beats `i` whenever `i < j`. -/
def revTour (n : Nat) : List Bool := (pairsLt n).map (fun _ => false)

/-- The 3-cycle: `0` beats `1`, `2` beats `0`, `1` beats `2`. -/
def cycle3 : List Bool := [true, false, true]

/-- (NC2) `winVec` genuinely reads the orientation and is not constant: reversing every game
of the 5-team witness reverses the win vector, and the 3-cycle tournament on 3 teams has win
vector `(1,1,1)` — repeated win counts, so nothing in the encoding forces distinctness (which
is exactly why the `k`-absoluteness hypothesis is needed at all). -/
theorem mse_2605396__ans2605536_nc2_winvec_is_not_constant :
    (natListEq (winVec 5 (revTour 5)) [0, 1, 2, 3, 4] = true)
  ∧ (natListEq (winVec 5 (revTour 5)) (winVec 5 (transTour 5)) = false)
  ∧ (natListEq (winVec 3 cycle3) [1, 1, 1] = true)
  ∧ (natSum (winVec 3 cycle3) = 3)
  ∧ (actualTotalWins 3 = 3) := by decide

/-- (NC3) `kAbsolute` / `hasTop` / `hasBot` are discriminating, not trivially true: the
3-cycle has neither a top nor a bottom element, so it is **not** `3`-absolute (whereas
`2`-absoluteness is automatic for every tournament, since one of any two teams beat the
other — which is why line `0005` insists on `k > 2`).  Hence Part 3's verification
that the 6-team witness *is* `3`-absolute carries real content. -/
theorem mse_2605396__ans2605536_nc3_kabsolute_discriminates :
    (hasTop 3 cycle3 [0, 1, 2] = false)
  ∧ (hasBot 3 cycle3 [0, 1, 2] = false)
  ∧ (kAbsolute 3 cycle3 3 = false)
  ∧ (kAbsolute 3 cycle3 2 = true)
  ∧ (hasTop 3 (transTour 3) [0, 1, 2] = true)
  ∧ (hasBot 3 (transTour 3) [0, 1, 2] = true)
  ∧ ((subsetsOfSize 3 (List.range 6)).length = 20) := by decide

/-- (NC4) `beatsB` really is a tournament relation on the witness: for every ordered pair of
distinct teams exactly one of `beatsB i j`, `beatsB j i` holds, and no team beats itself.  A
lookup that always returned `false` (which would make `hasTop` false and `kAbsolute` vacuous)
is ruled out. -/
theorem mse_2605396__ans2605536_nc4_beats_is_a_tournament :
    ((List.range 6).all (fun i => (List.range 6).all (fun j =>
        if i == j then !beatsB 6 (transTour 6) i j
        else (beatsB 6 (transTour 6) i j) != (beatsB 6 (transTour 6) j i))) = true)
  ∧ (beatsB 5 (transTour 5) 0 4 = true)
  ∧ (beatsB 5 (transTour 5) 4 0 = false)
  ∧ (beatsB 5 (revTour 5) 0 4 = false)
  ∧ (beatsB 5 (revTour 5) 4 0 = true) := by decide

/-- (NC5) The unattainable value `n` is not unattainable *in the encoding*: `hasWinCount` and
`realizesSet` happily return `true` for `5` on a hand-written vector containing `5`, and the
claimed hole multiset `{1,…,5}` is matched exactly by the (non-tournament) vector
`[1,2,3,4,5]`.  So the sweep's failures are facts about tournaments, not about a predicate
that can never fire. -/
theorem mse_2605396__ans2605536_nc5_claimed_values_are_expressible :
    (hasWinCount [5, 3, 2, 1, 0] 5 = true)
  ∧ (realizesSet [5, 4, 3, 2, 1] (claimedTopPeel 2 5) = true)
  ∧ (realizesSet [1, 2, 3, 4, 5] (claimedHoles 5) = true)
  ∧ (natListEq (sortNat [5, 4, 3, 2, 1]) (claimedHoles 5) = true)
  ∧ (natSum [1, 2, 3, 4, 5] = claimedTotalWins 5)
  ∧ (sweepOkWv 5 [1, 2, 3, 4, 5] = false) := by decide

/-- (NC6) The sweep predicate is not a tautology: `sweepOkWv` returns `false` on the
hypothetical win vector the proof describes (`{1,…,5}`, previous theorem) and `true` on the
real witness — the same five checks separate the two. -/
theorem mse_2605396__ans2605536_nc6_sweep_predicate_discriminates :
    (sweepOk 5 (transTour 5) = true)
  ∧ (sweepOk 5 (revTour 5) = true)
  ∧ (sweepOkWv 5 [1, 2, 3, 4, 5] = false)
  ∧ (sweepOkWv 5 [5, 3, 2, 1, 0] = false) := by decide

end mse_2605396__ans2605536

\end{Vlean}
\subsection*{Proof-safety information}
The automated scan found no unfinished proofs, new axiom declarations, unsafe declarations, external code hooks, or uses of Lean's compiled decision procedure. This scan concerns proof construction only; it does not establish that the theorem matches the source claim.
\textit{Technical scan counts:} native\_decide=0, sorry=0, axiom\_decl=0, unsafe=0, extern=0. \textit{Technical tier:} KERNEL-ONLY.
\textbf{A YES means:} Lean represents the same objects and assumptions, uses the same proposed counterexample, and the final theorem statement asserts the same contradiction. Successful compilation by itself is not enough.
\clearpage
\section*{Human sign-off}
Complete this block from your own review of Questions 1--3. Do not use the automated appendix below as an answer key.\par
\textbf{Reviewer name:} \rule{9cm}{0.4pt}\par
\textbf{Date:} \rule{4cm}{0.4pt}\par
\medskip\textbf{1. Does the proof verification document accurately represent the source and challenged claim?}\par
\answerchoices\par
\decisionexplanation
\medskip\textbf{2. Does the proposed counterexample satisfy the assumptions and make that claim false?}\par
\answerchoices\par
\decisionexplanation
\medskip\textbf{3. Does the final Lean theorem statement represent and verify that same counterexample?}\par
\answerchoices\par
\decisionexplanation
\medskip\textbf{Overall decision:}\par
\checkbox\ VALIDATED (all three answers are YES)\par
\checkbox\ REJECTED (at least one answer is NO)\par
\checkbox\ NEEDS ANOTHER REVIEWER (at least one CANNOT DETERMINE and no NO)\par
\medskip\textbf{Additional comments.} Use this space for any extra source, counterexample, calculation, or Lean details that support your decisions.\par\rule{\linewidth}{0.4pt}\par\vspace{7pt}\rule{\linewidth}{0.4pt}
\clearpage
\section*{Appendix --- Prior automated review (background only)}
\textbf{These results are not an answer key.} They come from the earlier Opus/Fable verification campaign and are printed only so an auditor can see what was previously checked. The human decisions above must be made independently. Raw status codes are retained in small type for traceability.
\textbf{Overall automated disposition:} Cleared after a follow-up reviewer resolved an automated concern.
\par\textit{Audit code:} \nolinkurl{PASS_ADJUDICATED}
\subsection*{Independent search with no access to the stored example}
\textbf{Result:} The search checked its entire stated range and found no counterexample there.
\par\textit{Audit code:} \nolinkurl{NO_WITNESS_EXHAUSTIVE}
\par\textbf{Checker's explanation:}
\begin{Vsrc}
No counterexample to the claim as it is used in this problem. In stratum B2 - the parameter regime the problem inhabits (N=93 teams, s=19, 93 > 2*19-5=33) - every one of the 17544 hypothesis-satisfying tournaments on up to 7 vertices is transitive with score sequence exactly 0,1,...,N-1. So under the problem's hypothesis the realized win-counts are all distinct AND fill a contiguous block, which is exactly what the quoted step asserts (teams with k+1,k+2,...,n wins). (Shortened here; the complete note remains in the machine-check JSON.)
\end{Vsrc}
\par\textbf{Cases checked:}
\begin{Vsrc}
Reading A (claim read context-free, 'the set of realized win-counts of a tournament is a contiguous block'): ALL 33864 labeled tournaments on N=3,4,5,6 vertices (2^C(N,2) each). Reading B (claim as USED here, i.e. applied to a tournament obeying the local source/sink hypothesis): ALL labeled tournaments on N=3..7 vertices (2^C(N,2) each; (Shortened here; the complete note remains in the machine-check JSON.)
\end{Vsrc}
\subsection*{Independent re-check of the proposed counterexample}
\textbf{Result:} A separate checker recomputed the proposed counterexample and confirmed it.
\par\textit{Audit code:} \nolinkurl{CONFIRMED}
\subsection*{Independent review of the source-to-Lean correspondence}
\textbf{Result:} A separate reviewer judged that Lean proves the same failure using the same example.
\par\textit{Audit code:} \nolinkurl{MATCHES}
\par\textbf{Theorem summary from that reviewer:}
\begin{Vsrc}
The two main theorems are `mse_2605396__ans2605536_witness5_refutes` and `mse_2605396__ans2605536_witness6_refutes`, both closed Bool conjunctions discharged by kernel `decide`. Encoding: a tournament on n teams is a `List Bool` indexed by `pairsLt n` (all pairs (i,j), i<j, lexicographic), entry true meaning "i beat j"; `winVec n o` is the out-degree vector of teams 0..n-1; `claimedTopPeel k n = [k+1,..., n]` is the quoted sentence's asserted set; (Shortened here; the complete note remains in the machine-check JSON.)
\end{Vsrc}
\par\textbf{Difference noted by that reviewer:}
\begin{Vsrc}
None material; all three links hold. (1) Definitions. `pairsLt`/`winnersOf`/`winCount`/`winVec` encode exactly "orientation of K_n" and out-degree = games won: every Bool list of length C(n,2) is a legal round-robin and vice versa, so no well-formedness side condition can exclude a counterexample, and ties are impossible by construction (matching the problem's "only possible results for a match is win and lose"). (Shortened here; the complete note remains in the machine-check JSON.)
\end{Vsrc}
\subsection*{Follow-up review of an automated concern}
\textbf{Result:} Follow-up review found that the record was right and an earlier checker was mistaken.
\par\textit{Audit code:} \nolinkurl{CHECKER_ERROR}.\par\textit{Confidence:} \nolinkurl{high}.\par\textit{Recommended action:} \nolinkurl{NONE}.
\par\textbf{Conclusion of the follow-up review:}
\begin{Vsrc}
The surrounding solution text binds n to the number of teams (line 0005: "In your case, n, k are 93, 19"; line 0007: "n pigeons (teams) and n holes (number of wins)"). In an n-team round robin every team plays exactly n-1 games, so every win-count lies in {0,...,n-1} and the value n is unattainable; since the asserted set {k+1,...,n} always contains n, it is unrealizable for every n and every k. Independent confirmations: total wins = C(n,2) = n(n-1)/2 while the sentence's hole multiset {1,...,n} weighs n(n+1)/2, an excess of exactly n (for n=93: 4278 vs 4371, excess 93). (Shortened here; the complete note remains in the machine-check JSON.)
\end{Vsrc}
\textit{This appendix is only background. You do not need any separate machine-output files to answer the three questions.}
\end{document}